% SPDX-License-Identifier: CC-BY-4.0
% Copyright (c) 2026 Martin Schröder <info@swedishembedded.com>

\section{Introduction}\label{sec:intro}

A deployed local model does not know what its users' documents, tools and repositories
say, and what it is taught by hand is lost the next time it is replaced. The
\emph{self-learning flywheel} is the closed loop that Splinter, Brain and Sven implement
to address this: take a source, write tasks from it, have an agent attempt them, keep
only what code can verify, project the verified experience into training data, train a
LoRA adapter (and later consolidate it into the weights), measure on material the
learner cannot influence, and release the result only if it is measurably better without
being worse at what it already did. The released artifact is an ordinary Brain adapter
that runs on plain \code{brain serve}; the next turn starts from it.

\paragraph{The result this work is for, and its status.} The practical question is
operational: \emph{can a team take its own git repository and tune a local agentic model
that is measurably better at completing unseen tasks in it, repeatedly, without losing what it
had?} The registered test is a paired comparison on repository tasks reserved before any
rollout (\cref{sec:protocol,sec:coding-results}). \textbf{It did not show the improvement}: the
candidate registered before the query reached a mean pass@1 of 23.0 against 20.0 for its base on
the 46 qualification families (18 wins, 12 losses, 16 ties, $p=0.18$). What did separate is the
supervision: against the same corpus, recipe and seed with every action supervised instead of the
decisive ones, the candidate won 26 families and lost 11 ($p=0.010$), and the six-seed development
comparison of the candidate against its base had given $p=0.040$.\cl{QUAL1,PROTO6} The paper
establishes what is needed to measure such a result honestly and reports what was measured: the
architecture, a task and verification funnel, a training substrate that fits a 27B model on one card, a
release gate, the requirements that make the loop work, and the coding-task protocol with
its results.

\paragraph{Scope.} The paper asks two questions. \emph{What must each stage of the loop do,
and what does the code contain for it, so that a turn of the loop can add a capability without
losing one?} (\cref{sec:architecture,sec:training,sec:substrate,sec:requirements}.)
\emph{Can the same loop be pointed at a repository and its effect on unseen tasks be decided
by a registered test?} (\cref{sec:protocol,sec:coding-results}, with the recipe for another
repository in \cref{sec:replicate}.) It claims
no general improvement of any benchmark score and contains no plans: a mechanism the code does not contain is marked planned, and a measurement that
was not made is stated as not made.

\paragraph{How claims are verified.} Every statement about the system or a measurement
carries the id of a row in a claims ledger (\code{data/claims.csv}, rendered as grey
bracketed ids). A \textsc{structure} claim is a fact about the code, a command line, the registration file or the
history of a repository and is re-checked by \code{make verify} against the repositories at
pinned commits that a fresh clone can reach (\cref{sec:repro}); a \textsc{measured} claim is a
number computed by \code{scripts/build_data.py} from the measurements recorded in \code{data/measurements.json};
a \textsc{computed} claim is a number derived by that script from a stated formula; a \textsc{cited} claim states what a published work does and points
at its bibliography entry. A structure check shows that a definition, constant or message exists
in a file; it does not show that the code behaves accordingly, which only the experiments do.

\paragraph{Contributions.}
(i)~A description of the division of labour among the three projects and of its enforcement
(\cref{sec:architecture}), and of the intake, admission, verification and leakage-control funnel
with its rejection reasons (\cref{sec:intake,sec:admission,sec:verify,sec:dataset}).
(ii)~The training, release-gate, rollback and lineage machinery (\cref{sec:training,sec:gate}) and
what the training substrate of one card holds and was measured to do for a 27B model
(\cref{sec:substrate}).
(iii)~The requirements that make the loop work, each tied to its implementation or marked
planned (\cref{sec:requirements}).
(iv)~A coding-task protocol fixed before any run, and its results: a registered confirmatory
query that does not show the candidate better than its base, and a same-seed supervision contrast
that separates (\cref{sec:protocol,sec:coding-results}), with a replication recipe whose steps are
marked available or planned (\cref{sec:replicate}).
(v)~The loop without a teacher, on a 4B model learning a language's standard library from the raw
manual: what a small model needs from the loop, full fine-tuning within one card, and the measured
rounds and final probes of a model that writes its own tasks against one whose base writes them
(\cref{sec:documentation}).
